\fichatitulo{12 · CONTEXTO PARLAMENTAR}{Revisão de literatura · 2024}{Pesquisa sobre discurso parlamentar}
{\small SKUBIC, Jure; FIŠER, Darja. \textit{Parliamentary Discourse Research in Political Science: Literature Review}. ParlaCLARIN IV, 2024, p. 1--11. \href{https://aclanthology.org/2024.parlaclarin-1.1/}{Fonte e versão}.\par}

\campo{Problema e objetivo}
Quais métodos a ciência política emprega para estudar registros parlamentares, e como corpora e técnicas de análise textual podem contribuir?

\campo{Principal contribuição · síntese interpretativa}
Organiza trabalhos por abordagem metodológica e aproxima pesquisa política e linguística de corpus. Destaca documentação e metadados como condições para reutilizar acervos de maneira contextualizada (seções 3 e 5).

\campo{Método}
Seleciona 24 artigos de 2012--2022, com filtros de disciplina, relevância/citação e explicitação metodológica. Apresenta até dois exemplos por abordagem, dada a heterogeneidade das perguntas (seção 3).

\campo{Resultado principal}
No conjunto selecionado, 12 dos 24 artigos usam análise de conteúdo; os demais abrangem discurso, sentimento, temas, métodos mistos e text-as-data. É uma descrição da seleção, não estimativa da prevalência em toda a ciência política (seção 3.2).

\campo{Excertos-chave · seleção literal}
\begin{itemize}
\excerto{Corpus da revisão}{24 relevant papers}{seção 3.1, p. 2.}
\excerto{Método}{content analysis}{seção 3.2, p. 2.}
\excerto{Documentação}{clearly and in detail document their compilation decisions}{seção 5, p. 9.}
\end{itemize}

\campo{Limitações declaradas e análise crítica}
Fonte: recorte temporal e seleção por relevância/citação, com apresentação de exemplos representativos (seção 3). \textbf{Análise da ficha:} esses filtros restringem a cobertura. A revisão não testa RAG, nem demonstra que acesso automatizado gere transparência ou participação.

\campo{Aproveitamento na dissertação · proposta}
Referencial teórico, Informação legislativa, transparência e acervos parlamentares: justificar metadados, contexto e documentação do corpus. Diferenciar recuperação de informação de interpretação política.

\registro{Fonte consultada nesta preparação: PDF publicado; consulta focal de Introdução, seleção/contagem de métodos (seção 3) e Discussão/Conclusão (seção 5). Estudos individuais da seção 4 não fichados. Excertos em inglês. Chave: \texttt{skubicFiser2024discourseReview}.}
